\documentclass[11pt,a4paper]{article}
\usepackage[T1]{fontenc}
\usepackage[utf8]{inputenc}
\usepackage[brazil]{babel}
\usepackage[margin=2.2cm]{geometry}
\usepackage{booktabs,longtable,array}
\usepackage[hidelinks]{hyperref}
\usepackage{xurl}

\newcolumntype{L}[1]{>{\raggedright\arraybackslash}p{#1}}

\title{Relatório do Dashboard\\\large Análise de Dados de Enchentes do INCT SiMAI}
\author{Henrique Nunes, Henzo Vasconcellos, Jayme, Rodrigo Gagliardi\\Coleta, Preparação e Análise de Dados --- PUCRS 2026/2}
\date{24 de novembro de 2026}

\begin{document}
\maketitle

\noindent\textbf{Link público:} \url{https://guaiba-poa.streamlit.app/}\\
\textbf{Repositório:} \url{https://github.com/henriquernunes/TrabalhoFinal-Coleta}

\begin{quote}
Os números de pessoas e prédios mostrados no dashboard são \textbf{estimativas de uma simulação}, não previsões nem contagens oficiais.
\end{quote}

\section{Ponto de partida: a ferramenta atual do INCT}
Antes de desenhar o dashboard, a equipe usou a página ``Flood --- INCT SiMAI'' e anotou o que ajuda e o que atrapalha.

\paragraph{O que funciona}
\begin{itemize}
  \item A simulação cobre os 94 bairros de Porto Alegre, de 3,0 a 10,0~m, em passos de 0,1~m.
  \item A pergunta central (``se o rio chegar a X metros, o que acontece?'') é clara e útil para defesa civil, imprensa e moradores.
\end{itemize}

\paragraph{O que atrapalha}
\begin{itemize}
  \item A página roda em servidor próprio, com instabilidade e lentidão: às vezes demora a carregar ou não responde.
  \item Falta um resumo da cidade inteira no nível escolhido (total de pessoas, prédios e bairros).
  \item Não há como ver, de uma vez, como o impacto cresce com o nível, nem onde fica o recorde de 2024 (5,35~m).
  \item Não fica claro que a população afetada é uma estimativa proporcional aos prédios, nem quais são as limitações dos dados.
\end{itemize}

\paragraph{O que levamos para o dashboard.} Cálculo sobre um arquivo limpo e pequeno, carregado de uma vez (resposta imediata ao mover o controle); indicadores da cidade no topo; mapa dos bairros que muda com o nível; curva por nível com o recorde marcado, para a cidade e para cada bairro; aviso de limitações sempre visível.

\section{Estrutura e navegação}
O dashboard é um aplicativo Streamlit com uma barra fixa no topo e três telas:
\begin{itemize}
  \item \textbf{Simulação INCT} (Tela 1);
  \item \textbf{Enchente 2024} (Tela 2, mapa) --- é a página inicial;
  \item \textbf{Histórico} (Tela 3, série histórica da ANA).
\end{itemize}
O enunciado prevê duas telas; com a autorização da professora, a parte das bases externas foi dividida em duas (mapa e histórico), porque o mapa ocupa a tela inteira. Cada uma cumpre sozinha os requisitos mínimos das bases externas: base coletada por código, pelo menos 3 visualizações, 1 filtro interativo, ligação com a Tela 1 e fonte com link e data de coleta. Cada entrada tem um endereço próprio (\texttt{?tela=inct}, \texttt{?tela=enchente}, \texttt{?tela=enchente\&painel=historico}), o que permite compartilhar o link direto de uma tela.

\section{Tela 1 --- Simulação do INCT}
\textbf{Objetivo.} Responder rapidamente: se o Guaíba chegar a X metros, quais bairros e quantas pessoas são afetados? Os números vêm apenas do JSON do INCT, depois da limpeza descrita em \texttt{limpeza\_dados.pdf}; o mapa usa também os limites oficiais dos bairros, só para desenhá-los.

\begin{longtable}{L{3.2cm} L{6.2cm} L{6.2cm}}
\toprule
Elemento & O que mostra & Por que essa escolha\\
\midrule
\endhead
Controle deslizante & Nível de 3,0 a 10,0~m, passo de 0,1~m; começa em 5,4~m (5,35~m arredondado) & Atualiza todos os elementos da tela. Começar perto do recorde dá um cenário concreto logo na abertura\\
Indicadores & Pessoas afetadas, prédios afetados e bairros atingidos & Respondem à pergunta em três números, sem precisar ler gráfico\\
Mapa dos bairros & Os 94 bairros coloridos pela porcentagem de prédios afetados no nível escolhido; bairros sem prédio atingido em cinza; no hover, porcentagem, prédios afetados/total e pessoas estimadas & Mostra de uma vez \emph{onde} ficam os bairros atingidos, o que a lista não mostra. Escala fixa de 0 a 100\%, em tons de azul, para comparar níveis sem cores alarmistas\\
Ranking & Os 10 bairros mais afetados, por pessoas ou por prédios (escolha do usuário) & Pessoas e prédios ordenam os bairros de forma diferente; a tabela traz os dois números e a porcentagem\\
Barras horizontais & Porcentagem de prédios afetados em cada bairro atingido, com prédios afetados/total no hover & Barras horizontais comportam nomes longos de bairros; a escala fixa de 0 a 100\% permite comparar níveis diferentes. Todos os bairros atingidos aparecem, com rolagem\\
Curva da cidade & População afetada estimada em cada nível, com linha tracejada em 5,35~m e linha no nível escolhido & Mostra o crescimento do impacto e situa o nível escolhido em relação ao recorde de 2024\\
Curva de um bairro & Seletor com os 94 bairros (começa no Sarandi, o mais atingido a 5,4~m); curva da porcentagem de prédios afetados por nível, com o recorde e o nível escolhido marcados; frase com o nível em que o bairro começa a ser atingido e o impacto no nível escolhido & Responde à pergunta de quem mora ou trabalha em um bairro específico, que o total da cidade não responde\\
Aviso de limitações & Caixa de informação abaixo das curvas e fonte dos dados & Exigência do enunciado e cuidado ético\\
\bottomrule
\end{longtable}

\textbf{Exemplos de leitura:} a 3,0~m, a simulação estima 32.112 pessoas em 24 bairros; a 5,4~m, 139.179 pessoas em 36 bairros (Sarandi, Farrapos e Menino Deus no topo); a 10,0~m, 271.875 pessoas em 54 bairros. No seletor, o Sarandi começa a ter prédios atingidos a 3,0~m e, a 5,4~m, tem 24.391 pessoas estimadas e 47,32\% dos prédios atingidos; o Menino Deus começa a 3,3~m; 40 bairros, como o Bom Fim, não têm prédios atingidos até 10,0~m.

\textbf{Mapa dos bairros.} O mapa coroplético, opcional no enunciado, foi incluído nesta tela. Os limites dos bairros são os da camada \texttt{BAIRROS} do GIS da Prefeitura (SMAMUS), a mesma usada no mapa da Tela 2, com os nomes já ajustados à grafia do INCT (acentos e apóstrofos mudam entre as fontes). O script \texttt{scripts/extrai\_bairros.py} copia esses limites para \texttt{dados/processados/bairros.geojson} e confere que os 94 nomes são os mesmos do INCT. O mapa é desenhado em SVG, sem imagem de fundo, para funcionar também em computadores sem aceleração gráfica.

\section{Tela 2 --- Enchente de maio de 2024 (mapa)}
\textbf{Objetivo.} Responder: o que outras bases acrescentam à simulação? A tela usa a ANA (o que o rio de fato fez, hora a hora) e a Prefeitura (onde a água chega em cada cota) e liga as duas ao INCT pelo nível do Guaíba e pelos bairros.

\begin{longtable}{L{3.2cm} L{6.2cm} L{6.2cm}}
\toprule
Elemento & O que mostra & Por que essa escolha\\
\midrule
\endhead
Linha do tempo (filtro) & Controle de hora a hora da enchente sobre o hidrograma do nível medido, com botão Reproduzir & Permite acompanhar a cheia no tempo; o hidrograma mostra subida, pico e descida no próprio controle\\
Mancha no mapa & Área abaixo da cota do nível daquela hora (camada da Prefeitura/SMAMUS) & Mostra onde a água chega em cada nível medido\\
Régua & Régua limnimétrica com a cota de inundação (3,00~m) e o recorde (5,35~m) & Lembra a régua real do Cais Mauá e dá referência visual da altura\\
Painel de leitura & Data, hora, nível, estação de origem, pessoas estimadas pelo INCT e eventos da enchente (comportas, casas de bomba) & Liga a ANA à Tela 1: o nível de cada hora é arredondado ao décimo e consultado no INCT\\
Bairros no hover & Pessoas estimadas, \% dos prédios e escolas (cadastro INEP) na área simulada daquela cota & Liga pelos bairros; as escolas indicam serviços na mancha simulada, sem afirmar danos reais\\
Botões & Realçar a água, prédios em 3D e imagem da inundação de 07/05/2024 publicada pela Prefeitura & Ajudam a ler o mapa e a comparar a mancha simulada com uma imagem real\\
\bottomrule
\end{longtable}

\section{Tela 3 --- Histórico do Guaíba}
\textbf{Objetivo.} Responder: com que frequência o Guaíba chega a níveis de inundação, e o que cada cheia representaria segundo o INCT? Usa a série diária da ANA desde 2014.

\begin{longtable}{L{3.2cm} L{6.2cm} L{6.2cm}}
\toprule
Elemento & O que mostra & Por que essa escolha\\
\midrule
\endhead
Intervalo de anos (filtro) & Seleciona os anos de 2014 a 2026 & Filtro interativo exigido; permite focar nos anos recentes\\
Dias acima da cota & Barras com os dias por ano em que o nível passou da cota (3,00~m no Cais Mauá; 3,60~m no Gasômetro) & Responde à frequência das cheias: nenhum dia entre 2014 e 2022; 6 em 2023, 32 em 2024 e 2 em 2025\\
Maior nível do ano & Linha com o pico anual e linha tracejada em 5,35~m & Mostra o quanto 2024 se afasta dos outros anos\\
Pessoas estimadas no pico & Barras com a estimativa do INCT para o pico de cada ano (só acima de 3,0~m) & Liga a ANA à Tela 1 pelo nível: 56.027 pessoas em 2023, 139.179 em 2024 e 32.112 em 2025, todas como cenário hipotético\\
\bottomrule
\end{longtable}

\section{Fontes e datas de coleta}
Cada tela mostra a fonte com link e a data de coleta:
\begin{itemize}
  \item \textbf{ANA} --- web service de telemetria \url{https://telemetriaws1.ana.gov.br/ServiceANA.asmx}, estações \texttt{87450004} (Cais Mauá C6) e \texttt{87444000} (Usina do Gasômetro); coleta em 28/10/2026.
  \item \textbf{Prefeitura de Porto Alegre (SMAMUS/DMAE)} --- \url{https://gis-smamus.portoalegre.rs.gov.br/server/rest/services/Hosted}: mancha por cota, bairros, eventos de 2024 e escolas (cadastro INEP); coleta em 09/11/2026. Os limites dos bairros também desenham o mapa da Tela 1.
  \item \textbf{INCT SiMAI} --- JSON fornecido pela disciplina, limpo pelos scripts do repositório.
\end{itemize}

\section{Como usar}
\begin{enumerate}
  \item Abra o link público. A página inicial é o mapa de maio de 2024.
  \item No mapa, arraste a linha do tempo ou clique em \textbf{Reproduzir}. Passe o mouse sobre um bairro para ver a estimativa do INCT e as escolas na área simulada.
  \item Clique em \textbf{Histórico} e ajuste o intervalo de anos para ver frequência e picos das cheias.
  \item Clique em \textbf{Simulação INCT}, mova o controle do nível e leia indicadores, mapa dos bairros, ranking, barras e curva da cidade. Passe o mouse sobre o mapa para ver os números de cada bairro e troque o ranking entre pessoas e prédios.
  \item No fim da Tela 1, escolha um bairro no seletor para ver a curva dele e a partir de que nível ele começa a ser atingido.
\end{enumerate}
Para rodar localmente, siga o \texttt{README.md}: criar o ambiente, coletar (\texttt{scripts/coletar\_tudo.py}), limpar e processar (\texttt{limpa\_inct.py}, \texttt{processa\_ana.py}, \texttt{gera\_mapa.py}, \texttt{extrai\_bairros.py}) e executar \texttt{streamlit run app/app.py}.

\section{Limitações}
\begin{itemize}
  \item \textbf{Estimativas, não contagens.} A população afetada supõe moradores distribuídos igualmente entre todos os prédios, inclusive os não residenciais.
  \item \textbf{Só o nível do rio.} A simulação e a mancha por cota não consideram diques, comportas, casas de bomba, chuva local nem água retida na descida.
  \item \textbf{Correções no JSON.} Valores de prédios afetados que caíam com o rio subindo foram interpolados, e a porcentagem foi recalculada sem arredondar para cima (detalhes em \texttt{limpeza\_dados.pdf}).
  \item \textbf{Duas estações.} O Cais Mauá parou em 02/05/2024; de 03/05/2024 a 11/06/2025 a série usa o Gasômetro. As estações têm referenciais diferentes e não foram ajustadas entre si; a lacuna de 18 horas entre elas foi interpolada.
  \item \textbf{Cobertura desigual.} Há lacunas na telemetria (por exemplo, 2014 tem 126 dias com leitura); os dias sem dados não entram na contagem.
  \item \textbf{Arredondamento.} Para consultar o INCT, o nível é arredondado ao décimo (5,35~m vira 5,4~m); abaixo de 3,0~m não há estimativa.
  \item \textbf{Escolas.} Indicam pontos do cadastro dentro da mancha simulada, não escolas que de fato alagaram.
  \item \textbf{Mapa da Tela 1.} Cada bairro recebe uma única cor, pela porcentagem de prédios afetados no bairro todo; o mapa não mostra em que parte do bairro ficam esses prédios.
\end{itemize}

\section{O que mudou desde a Fase 3}
\begin{itemize}
  \item A régua do mapa passou a ser redesenhada quando a janela muda de tamanho; antes, as marcas ficavam na escala antiga e não batiam com o nível mostrado.
  \item O dashboard foi publicado no Streamlit Community Cloud, e o link foi incluído no README e neste relatório.
  \item O README passou a listar todos os relatórios e o uso de IA na Fase 3.
  \item O wireframe e este relatório passaram a tratar o Histórico como Tela 3, como ele aparece na navegação do dashboard.
  \item A Tela 1 ganhou o mapa dos bairros colorido pela porcentagem de prédios afetados, opcional no enunciado.
  \item A Tela 1 ganhou o seletor de bairro, com a curva do bairro por nível e o nível em que ele começa a ser atingido.
  \item Os testes automatizados foram ampliados de 5 para 51. Eles cobrem a limpeza do INCT e da ANA, o cruzamento entre as bases, as funções do mapa e as telas, e conferem que os scripts reproduzem exatamente os arquivos processados do repositório e os números do relatório de limpeza.
\end{itemize}

\section{Uso de IA}
O Claude Code (Anthropic) foi usado como apoio para escrever o código do mapa e das telas em Streamlit (incluindo o mapa dos bairros e o seletor de bairro da Tela 1), os testes, o wireframe e o rascunho deste relatório. A equipe revisou o código e o texto e consegue explicar cada parte. O registro completo está na tabela ``Uso de IA'' do README.

\end{document}
